% Parte do material do minicurso "Escrita de Artigos Científicos com LaTeX" (Conecta Hub)
% Exemplo 04 – Tabelas e figuras
% Usa apenas recursos nativos (sem imagens externas), para compilar sem arquivos extras.

\documentclass[12pt, a4paper]{article}
\usepackage[utf8]{inputenc}
\usepackage[T1]{fontenc}
\usepackage[brazil]{babel}
\usepackage{booktabs}                    % linhas de tabela profissionais (\toprule, \midrule, \bottomrule)
\usepackage{graphicx}                    % inclusão de imagens (\includegraphics)
\usepackage{tikz}                        % desenho vetorial simples

\title{Tabelas e Figuras}
\author{Seu Nome}
\date{}

\begin{document}

\maketitle

\section{Tabela}

\begin{table}[h]
    \centering
    \caption{Desempenho de três algoritmos.}
    \label{tab:desempenho}
    \begin{tabular}{lcc}
        \toprule
        Algoritmo & Tempo (s) & Memória (MB) \\
        \midrule
        Alfa      & 1,20      & 35           \\
        Beta      & 0,85      & 42           \\
        Gama      & 2,10      & 28           \\
        \bottomrule
    \end{tabular}
\end{table}

Na Tabela~\ref{tab:desempenho}, o algoritmo Beta é o mais rápido.

\section{Figura}

\begin{figure}[h]
    \centering
    \begin{tikzpicture}
        \draw[->] (0,0) -- (4,0) node[right] {$x$};
        \draw[->] (0,0) -- (0,3) node[above] {$y$};
        \draw[thick, blue] (0,0) .. controls (1,2.5) .. (2.5,0.5) -- (3.5,2);
    \end{tikzpicture}
    \caption{Exemplo de figura desenhada com TikZ.}
    \label{fig:exemplo}
\end{figure}

A Figura~\ref{fig:exemplo} mostra um desenho vetorial.

\medskip
Para usar uma imagem sua (PNG/JPG/PDF), coloque o arquivo na pasta e escreva:

\begin{verbatim}
\includegraphics[width=0.6\textwidth]{minha_imagem.png}
\end{verbatim}

\end{document}
